\newcommand{\coursemodule}{Module 10}
% ════════════════════════════════════════════════════════════════════════════
% Pixels, Frames & Transformers: shared Beamer preamble (WUR style)
% Adapted from the preamble of Bags, Vectors & Transformers.
%
% Use in a module deck (from <part>/<module>/slides/):
%   \newcommand{\coursemodule}{Module 01}
%   % ════════════════════════════════════════════════════════════════════════════
% Pixels, Frames & Transformers: shared Beamer preamble (WUR style)
% Adapted from the preamble of Bags, Vectors & Transformers.
%
% Use in a module deck (from <part>/<module>/slides/):
%   \newcommand{\coursemodule}{Module 01}
%   % ════════════════════════════════════════════════════════════════════════════
% Pixels, Frames & Transformers: shared Beamer preamble (WUR style)
% Adapted from the preamble of Bags, Vectors & Transformers.
%
% Use in a module deck (from <part>/<module>/slides/):
%   \newcommand{\coursemodule}{Module 01}
%   \input{../../../shared/preamble.tex}
%   \begin{document}
%   \moduletitle{Images as Data}
%   ...
% Compile with: latexmk -pdf <deck>.tex
% ════════════════════════════════════════════════════════════════════════════

\documentclass[aspectratio=169]{beamer}

% ── Packages ────────────────────────────────────────────────────────────────
\usepackage[T1]{fontenc}
\usepackage[utf8]{inputenc}
\usepackage{lmodern}
\usepackage{microtype}
\usepackage{xcolor}
\usepackage{tikz}
\usetikzlibrary{arrows.meta, positioning, calc, shapes.geometric, decorations.pathreplacing, matrix, fit}
\usepackage{booktabs}
\usepackage{fontawesome5}
\usepackage{array}
\usepackage[most]{tcolorbox}

\providecommand{\coursemodule}{}

% ── WUR color palette ───────────────────────────────────────────────────────
\definecolor{wurgreen}{HTML}{34B233}
\definecolor{darkgreen}{HTML}{1F6B1E}
\definecolor{lightgreen}{HTML}{E8F7E8}
\definecolor{darknavy}{HTML}{1A1A2E}
\definecolor{midgray}{HTML}{6B7280}
\definecolor{lightgray}{HTML}{F3F4F6}
\definecolor{bodytext}{HTML}{1F2937}
\definecolor{codegreen}{HTML}{9AE6B4}
\definecolor{warnred}{HTML}{FEE2E2}
\definecolor{warnborder}{HTML}{EF4444}

% ── Beamer theme ────────────────────────────────────────────────────────────
\usetheme{default}
\usecolortheme{default}
\usefonttheme{professionalfonts}
\setbeamertemplate{navigation symbols}{}
\setbeamertemplate{itemize items}{\textcolor{wurgreen}{\textbullet}}
\setbeamertemplate{itemize subitem}{\textcolor{midgray}{--}}
\setbeamertemplate{enumerate items}{\textcolor{wurgreen}{\textbf{\insertenumlabel.}}}
\setbeamercolor{frametitle}{fg=white, bg=darknavy}
\setbeamertemplate{frametitle}{%
  \begin{beamercolorbox}[wd=\paperwidth, ht=1.1cm, dp=0.2cm, leftskip=0.6cm]{frametitle}
    \usebeamerfont{frametitle}\insertframetitle
  \end{beamercolorbox}%
}
\setbeamerfont{frametitle}{size=\large, series=\bfseries}
\setbeamercolor{background canvas}{bg=white}
\setbeamercolor{normal text}{fg=bodytext}
\setbeamerfont{normal text}{size=\small}
\setbeamertemplate{footline}{%
  \begin{beamercolorbox}[wd=\paperwidth, ht=0.4cm, dp=0.15cm]{footline}%
    \hspace{0.5cm}%
    \textcolor{midgray}{\tiny Pixels, Frames \& Transformers · \coursemodule\ · Wageningen University \& Research}%
    \hfill%
    \textcolor{midgray}{\tiny \insertframenumber}%
    \hspace{0.4cm}%
  \end{beamercolorbox}%
}

% ── Boxes ───────────────────────────────────────────────────────────────────
\tcbset{wurbase/.style={arc=3pt, boxrule=0.6pt, left=7pt, right=7pt, top=5pt, bottom=5pt,
                        fontupper=\small, before skip=4pt, after skip=4pt}}
\newtcolorbox{greenbox}[1][]{wurbase, colback=lightgreen, colframe=wurgreen, #1}
\newtcolorbox{graybox}[1][]{wurbase, colback=lightgray, colframe=lightgray, #1}
\newtcolorbox{navybox}[1][]{wurbase, colback=darknavy, colframe=darknavy, colupper=white, #1}
\newtcolorbox{warnbox}[1][]{wurbase, colback=warnred, colframe=warnborder, #1}
\newtcolorbox{codebox}[1][]{wurbase, colback=darknavy, colframe=darknavy, colupper=white,
                            fontupper=\ttfamily\scriptsize, left=6pt, top=4pt, bottom=4pt, #1}

% ── Title slide ─────────────────────────────────────────────────────────────
\newcommand{\moduletitle}[1]{%
{
\setbeamertemplate{footline}{}
\begin{frame}[plain]
  \begin{tikzpicture}[remember picture, overlay]
    \fill[darknavy] (current page.south west) rectangle (current page.north east);
    \fill[wurgreen] (current page.south west) rectangle ([xshift=0.45cm]current page.north west);
    \node[anchor=west, text=white, font=\huge\bfseries, text width=13cm]
      at ([xshift=1.1cm, yshift=1.2cm]current page.west) {Pixels, Frames \& Transformers};
    \node[anchor=west, text=wurgreen, font=\large, text width=13cm]
      at ([xshift=1.1cm, yshift=0.05cm]current page.west) {#1 · \coursemodule};
    \node[anchor=west, text=midgray, font=\small, text width=12cm]
      at ([xshift=1.1cm, yshift=-1.5cm]current page.west)
      {Denise J. Roth, PhD Candidate\\
       Strategic Communication Group · Wageningen University \& Research};
  \end{tikzpicture}
\end{frame}
}}

% ── Small helpers ───────────────────────────────────────────────────────────
\newcommand{\keyterm}[1]{\textbf{\textcolor{darkgreen}{#1}}}
\newcommand{\muted}[1]{\textcolor{midgray}{#1}}
\newcommand{\credit}[1]{{\tiny\textcolor{midgray}{#1}}}
\newcommand{\notebookslide}[1]{%
\begin{frame}{In the Notebook}
\vspace{0.1cm}
#1
\end{frame}}
\newcommand{\nbitem}[2]{\item[\textcolor{wurgreen}{#1}] #2}

%   \begin{document}
%   \moduletitle{Images as Data}
%   ...
% Compile with: latexmk -pdf <deck>.tex
% ════════════════════════════════════════════════════════════════════════════

\documentclass[aspectratio=169]{beamer}

% ── Packages ────────────────────────────────────────────────────────────────
\usepackage[T1]{fontenc}
\usepackage[utf8]{inputenc}
\usepackage{lmodern}
\usepackage{microtype}
\usepackage{xcolor}
\usepackage{tikz}
\usetikzlibrary{arrows.meta, positioning, calc, shapes.geometric, decorations.pathreplacing, matrix, fit}
\usepackage{booktabs}
\usepackage{fontawesome5}
\usepackage{array}
\usepackage[most]{tcolorbox}

\providecommand{\coursemodule}{}

% ── WUR color palette ───────────────────────────────────────────────────────
\definecolor{wurgreen}{HTML}{34B233}
\definecolor{darkgreen}{HTML}{1F6B1E}
\definecolor{lightgreen}{HTML}{E8F7E8}
\definecolor{darknavy}{HTML}{1A1A2E}
\definecolor{midgray}{HTML}{6B7280}
\definecolor{lightgray}{HTML}{F3F4F6}
\definecolor{bodytext}{HTML}{1F2937}
\definecolor{codegreen}{HTML}{9AE6B4}
\definecolor{warnred}{HTML}{FEE2E2}
\definecolor{warnborder}{HTML}{EF4444}

% ── Beamer theme ────────────────────────────────────────────────────────────
\usetheme{default}
\usecolortheme{default}
\usefonttheme{professionalfonts}
\setbeamertemplate{navigation symbols}{}
\setbeamertemplate{itemize items}{\textcolor{wurgreen}{\textbullet}}
\setbeamertemplate{itemize subitem}{\textcolor{midgray}{--}}
\setbeamertemplate{enumerate items}{\textcolor{wurgreen}{\textbf{\insertenumlabel.}}}
\setbeamercolor{frametitle}{fg=white, bg=darknavy}
\setbeamertemplate{frametitle}{%
  \begin{beamercolorbox}[wd=\paperwidth, ht=1.1cm, dp=0.2cm, leftskip=0.6cm]{frametitle}
    \usebeamerfont{frametitle}\insertframetitle
  \end{beamercolorbox}%
}
\setbeamerfont{frametitle}{size=\large, series=\bfseries}
\setbeamercolor{background canvas}{bg=white}
\setbeamercolor{normal text}{fg=bodytext}
\setbeamerfont{normal text}{size=\small}
\setbeamertemplate{footline}{%
  \begin{beamercolorbox}[wd=\paperwidth, ht=0.4cm, dp=0.15cm]{footline}%
    \hspace{0.5cm}%
    \textcolor{midgray}{\tiny Pixels, Frames \& Transformers · \coursemodule\ · Wageningen University \& Research}%
    \hfill%
    \textcolor{midgray}{\tiny \insertframenumber}%
    \hspace{0.4cm}%
  \end{beamercolorbox}%
}

% ── Boxes ───────────────────────────────────────────────────────────────────
\tcbset{wurbase/.style={arc=3pt, boxrule=0.6pt, left=7pt, right=7pt, top=5pt, bottom=5pt,
                        fontupper=\small, before skip=4pt, after skip=4pt}}
\newtcolorbox{greenbox}[1][]{wurbase, colback=lightgreen, colframe=wurgreen, #1}
\newtcolorbox{graybox}[1][]{wurbase, colback=lightgray, colframe=lightgray, #1}
\newtcolorbox{navybox}[1][]{wurbase, colback=darknavy, colframe=darknavy, colupper=white, #1}
\newtcolorbox{warnbox}[1][]{wurbase, colback=warnred, colframe=warnborder, #1}
\newtcolorbox{codebox}[1][]{wurbase, colback=darknavy, colframe=darknavy, colupper=white,
                            fontupper=\ttfamily\scriptsize, left=6pt, top=4pt, bottom=4pt, #1}

% ── Title slide ─────────────────────────────────────────────────────────────
\newcommand{\moduletitle}[1]{%
{
\setbeamertemplate{footline}{}
\begin{frame}[plain]
  \begin{tikzpicture}[remember picture, overlay]
    \fill[darknavy] (current page.south west) rectangle (current page.north east);
    \fill[wurgreen] (current page.south west) rectangle ([xshift=0.45cm]current page.north west);
    \node[anchor=west, text=white, font=\huge\bfseries, text width=13cm]
      at ([xshift=1.1cm, yshift=1.2cm]current page.west) {Pixels, Frames \& Transformers};
    \node[anchor=west, text=wurgreen, font=\large, text width=13cm]
      at ([xshift=1.1cm, yshift=0.05cm]current page.west) {#1 · \coursemodule};
    \node[anchor=west, text=midgray, font=\small, text width=12cm]
      at ([xshift=1.1cm, yshift=-1.5cm]current page.west)
      {Denise J. Roth, PhD Candidate\\
       Strategic Communication Group · Wageningen University \& Research};
  \end{tikzpicture}
\end{frame}
}}

% ── Small helpers ───────────────────────────────────────────────────────────
\newcommand{\keyterm}[1]{\textbf{\textcolor{darkgreen}{#1}}}
\newcommand{\muted}[1]{\textcolor{midgray}{#1}}
\newcommand{\credit}[1]{{\tiny\textcolor{midgray}{#1}}}
\newcommand{\notebookslide}[1]{%
\begin{frame}{In the Notebook}
\vspace{0.1cm}
#1
\end{frame}}
\newcommand{\nbitem}[2]{\item[\textcolor{wurgreen}{#1}] #2}

%   \begin{document}
%   \moduletitle{Images as Data}
%   ...
% Compile with: latexmk -pdf <deck>.tex
% ════════════════════════════════════════════════════════════════════════════

\documentclass[aspectratio=169]{beamer}

% ── Packages ────────────────────────────────────────────────────────────────
\usepackage[T1]{fontenc}
\usepackage[utf8]{inputenc}
\usepackage{lmodern}
\usepackage{microtype}
\usepackage{xcolor}
\usepackage{tikz}
\usetikzlibrary{arrows.meta, positioning, calc, shapes.geometric, decorations.pathreplacing, matrix, fit}
\usepackage{booktabs}
\usepackage{fontawesome5}
\usepackage{array}
\usepackage[most]{tcolorbox}

\providecommand{\coursemodule}{}

% ── WUR color palette ───────────────────────────────────────────────────────
\definecolor{wurgreen}{HTML}{34B233}
\definecolor{darkgreen}{HTML}{1F6B1E}
\definecolor{lightgreen}{HTML}{E8F7E8}
\definecolor{darknavy}{HTML}{1A1A2E}
\definecolor{midgray}{HTML}{6B7280}
\definecolor{lightgray}{HTML}{F3F4F6}
\definecolor{bodytext}{HTML}{1F2937}
\definecolor{codegreen}{HTML}{9AE6B4}
\definecolor{warnred}{HTML}{FEE2E2}
\definecolor{warnborder}{HTML}{EF4444}

% ── Beamer theme ────────────────────────────────────────────────────────────
\usetheme{default}
\usecolortheme{default}
\usefonttheme{professionalfonts}
\setbeamertemplate{navigation symbols}{}
\setbeamertemplate{itemize items}{\textcolor{wurgreen}{\textbullet}}
\setbeamertemplate{itemize subitem}{\textcolor{midgray}{--}}
\setbeamertemplate{enumerate items}{\textcolor{wurgreen}{\textbf{\insertenumlabel.}}}
\setbeamercolor{frametitle}{fg=white, bg=darknavy}
\setbeamertemplate{frametitle}{%
  \begin{beamercolorbox}[wd=\paperwidth, ht=1.1cm, dp=0.2cm, leftskip=0.6cm]{frametitle}
    \usebeamerfont{frametitle}\insertframetitle
  \end{beamercolorbox}%
}
\setbeamerfont{frametitle}{size=\large, series=\bfseries}
\setbeamercolor{background canvas}{bg=white}
\setbeamercolor{normal text}{fg=bodytext}
\setbeamerfont{normal text}{size=\small}
\setbeamertemplate{footline}{%
  \begin{beamercolorbox}[wd=\paperwidth, ht=0.4cm, dp=0.15cm]{footline}%
    \hspace{0.5cm}%
    \textcolor{midgray}{\tiny Pixels, Frames \& Transformers · \coursemodule\ · Wageningen University \& Research}%
    \hfill%
    \textcolor{midgray}{\tiny \insertframenumber}%
    \hspace{0.4cm}%
  \end{beamercolorbox}%
}

% ── Boxes ───────────────────────────────────────────────────────────────────
\tcbset{wurbase/.style={arc=3pt, boxrule=0.6pt, left=7pt, right=7pt, top=5pt, bottom=5pt,
                        fontupper=\small, before skip=4pt, after skip=4pt}}
\newtcolorbox{greenbox}[1][]{wurbase, colback=lightgreen, colframe=wurgreen, #1}
\newtcolorbox{graybox}[1][]{wurbase, colback=lightgray, colframe=lightgray, #1}
\newtcolorbox{navybox}[1][]{wurbase, colback=darknavy, colframe=darknavy, colupper=white, #1}
\newtcolorbox{warnbox}[1][]{wurbase, colback=warnred, colframe=warnborder, #1}
\newtcolorbox{codebox}[1][]{wurbase, colback=darknavy, colframe=darknavy, colupper=white,
                            fontupper=\ttfamily\scriptsize, left=6pt, top=4pt, bottom=4pt, #1}

% ── Title slide ─────────────────────────────────────────────────────────────
\newcommand{\moduletitle}[1]{%
{
\setbeamertemplate{footline}{}
\begin{frame}[plain]
  \begin{tikzpicture}[remember picture, overlay]
    \fill[darknavy] (current page.south west) rectangle (current page.north east);
    \fill[wurgreen] (current page.south west) rectangle ([xshift=0.45cm]current page.north west);
    \node[anchor=west, text=white, font=\huge\bfseries, text width=13cm]
      at ([xshift=1.1cm, yshift=1.2cm]current page.west) {Pixels, Frames \& Transformers};
    \node[anchor=west, text=wurgreen, font=\large, text width=13cm]
      at ([xshift=1.1cm, yshift=0.05cm]current page.west) {#1 · \coursemodule};
    \node[anchor=west, text=midgray, font=\small, text width=12cm]
      at ([xshift=1.1cm, yshift=-1.5cm]current page.west)
      {Denise J. Roth, PhD Candidate\\
       Strategic Communication Group · Wageningen University \& Research};
  \end{tikzpicture}
\end{frame}
}}

% ── Small helpers ───────────────────────────────────────────────────────────
\newcommand{\keyterm}[1]{\textbf{\textcolor{darkgreen}{#1}}}
\newcommand{\muted}[1]{\textcolor{midgray}{#1}}
\newcommand{\credit}[1]{{\tiny\textcolor{midgray}{#1}}}
\newcommand{\notebookslide}[1]{%
\begin{frame}{In the Notebook}
\vspace{0.1cm}
#1
\end{frame}}
\newcommand{\nbitem}[2]{\item[\textcolor{wurgreen}{#1}] #2}


\begin{document}

\moduletitle{Vision-Language Models as Annotators}

% ─────────────────────────────────────────────────────────────────────────────
\begin{frame}{A Research Assistant That Can See}
\begin{columns}[T]
\begin{column}{0.5\textwidth}
A \keyterm{vision-language model} (VLM) is a language model that can look at images. You can:

\begin{itemize}
\item ask what a picture shows
\item have it fill in a coding sheet
\item have it describe a scene for a text pipeline
\end{itemize}

\vspace{0.15cm}
Small open VLMs now run on a free Colab GPU.
\end{column}
\begin{column}{0.46\textwidth}
\begin{greenbox}
\textbf{Treat it like a new research assistant}\\
Give it your codebook, check its work against yours, study its mistakes, and trust it only where
it has earned that trust.
\end{greenbox}
\muted{\footnotesize The same workflow as Day 4 of Bags, Vectors \& Transformers: prompt,
validate, distill.}
\end{column}
\end{columns}
\end{frame}

% ─────────────────────────────────────────────────────────────────────────────
\begin{frame}{How a VLM Is Built}
\vspace{0.1cm}
\centering
\begin{tikzpicture}[
  b/.style={draw=midgray, fill=lightgray, rounded corners=3pt, font=\footnotesize, align=center, inner sep=5pt},
  g/.style={draw=wurgreen, fill=lightgreen, rounded corners=3pt, font=\footnotesize, align=center, inner sep=5pt},
  n/.style={draw=darknavy, fill=darknavy, text=white, rounded corners=3pt, font=\footnotesize, align=center, inner sep=5pt},
  arr/.style={-{Stealth[length=2mm]}, thick, color=midgray}]
\node[b, text width=1.6cm] (img) at (0, 0.9) {image};
\node[g, text width=2.5cm] (enc) at (2.75, 0.9) {vision encoder\\ \muted{a ViT, module 08}};
\node[g, text width=2.5cm] (con) at (5.75, 0.9) {connector\\ \muted{into the LM's space}};
\node[b, text width=2.5cm] (txt) at (5.75, -0.7) {your prompt};
\node[n, text width=2.7cm] (lm) at (8.85, 0.1) {language model\\ reads image tokens\\ and the prompt};
\node[b, text width=1.9cm] (ans) at (11.6, 0.1) {answer,\\ token by token};
\draw[arr] (img) -- (enc);
\draw[arr] (enc) -- (con);
\draw[arr] (con.east) -- (lm.north west);
\draw[arr] (txt.east) -- (lm.south west);
\draw[arr] (lm) -- (ans);
\end{tikzpicture}

\vspace{0.35cm}
\raggedright
Qwen3-VL merges each 2 $\times$ 2 block of patches into one token, so a 768-pixel image becomes a
few hundred tokens. The \textbf{language model} is where the knowledge, the instruction following
and the hallucinations come from.
\end{frame}

% ─────────────────────────────────────────────────────────────────────────────
\begin{frame}{Open Models That Fit a Free GPU}
\begin{columns}[T]
\begin{column}{0.58\textwidth}
{\footnotesize
\begin{tabular}{@{}>{\raggedright\arraybackslash}p{2.6cm}>{\raggedright\arraybackslash}p{1.6cm}>{\raggedright\arraybackslash}p{3.0cm}@{}}
\toprule
model & license & note \\
\midrule
\textbf{Qwen3-VL} 2B/4B & Apache 2.0 & our default, runs on a T4 \\
Qwen3.5 2B/4B & Apache 2.0 & newer, but needs a GPU with bfloat16 (not a T4) \\
Molmo2 4B & Apache 2.0 & its training data is open too \\
\bottomrule
\end{tabular}}

\vspace{0.2cm}
\muted{\footnotesize On a T4 (no bfloat16) the notebook runs the 2B model in float32: about 9 GB,
slower, but float16 can overflow and garble answers.}
\end{column}
\begin{column}{0.38\textwidth}
\begin{graybox}
\textbf{Why open weights matter}\\
A versioned open model gives the same answers next year. A commercial API can change the model
behind it without notice, and you can't inspect what it was trained on.
\end{graybox}
\muted{\footnotesize This landscape changes every few months; check before you start a project.}
\end{column}
\end{columns}
\end{frame}

% ─────────────────────────────────────────────────────────────────────────────
\begin{frame}[fragile]{Your Codebook as a Prompt}
\begin{columns}[T]
\begin{column}{0.52\textwidth}
\begin{codebox}
{\scriptsize\ttfamily
You are coding photos for a study of protest.\\
Code only what is visible in the photo.\\[3pt]
- protest: people engaged in a protest: \ldots\\
\hspace*{1em}Code 0 for \ldots a festival, a market \ldots\\
- crowd: roughly ten or more people \ldots\\
- police: \ldots Security guards don't count.\\
- signs: \ldots Street names don't count.\\[3pt]
Answer with JSON only, in exactly this format:\\
\{"protest": 0, "crowd": 0, "police": 0, "signs": 0\}}
\end{codebox}
\end{column}
\begin{column}{0.44\textwidth}
\begin{itemize}
\item Definitions \emph{and} rules for hard cases, as in module 04
\item \keyterm{JSON} output: easy to read into a table, and obvious when the format breaks
\item \keyterm{Greedy decoding}: same image, same prompt, same answer
\item Keep every raw answer, and count the ones that failed
\end{itemize}
\end{column}
\end{columns}
\end{frame}

% ─────────────────────────────────────────────────────────────────────────────
\begin{frame}{Validate It Like a Human Coder}
\begin{columns}[T]
\begin{column}{0.5\textwidth}
Per variable, against your codes:
\begin{itemize}
\item agreement and \keyterm{kappa} (rare variables can have high agreement and low kappa)
\item the \textbf{share coded yes}, by you and by the model: a model that says yes too often
      biases every prevalence you estimate
\item errors per group or context, if you compare groups (module 07)
\end{itemize}
\end{column}
\begin{column}{0.46\textwidth}
\begin{greenbox}
\textbf{Look at the disagreements}\\
Some are model errors, some are cases the codebook doesn't settle, and some are your own slips.
Human codes are the benchmark, not the truth.
\end{greenbox}
\begin{warnbox}
Answers that failed are often the hard images. Never drop them silently.
\end{warnbox}
\end{column}
\end{columns}
\end{frame}

% ─────────────────────────────────────────────────────────────────────────────
\begin{frame}{Hallucination}
\begin{columns}[T]
\begin{column}{0.5\textwidth}
A language model writes the most plausible continuation, and a plausible description of a protest
includes police and signs, whether or not the photo shows them \muted{(Rohrbach et al., 2018)}.

\vspace{0.2cm}
\textbf{A quick test}: on images you coded \emph{without} police, compare
\begin{itemize}
\item a neutral question: ``Is there a police officer? Answer yes or no.''
\item a leading one: ``Describe what the police officers are doing.''
\end{itemize}
\end{column}
\begin{column}{0.46\textwidth}
\begin{warnbox}
\textbf{Leading questions invite invention.} Small models often describe officers who aren't
there, even after answering ``no'' to the neutral question.
\end{warnbox}
\begin{greenbox}
Ask neutral questions with an explicit ``0 = not visible'', and treat open descriptions as claims
to check.
\end{greenbox}
\end{column}
\end{columns}
\end{frame}

% ─────────────────────────────────────────────────────────────────────────────
\begin{frame}{Sensitivity Checks}
\begin{columns}[T]
\begin{column}{0.48\textwidth}
\begin{graybox}[height=3.4cm]
\textbf{Prompt wording}\\
A shorter prompt with the same variables in another order. Report how often the two agree, as you
would for two human coders.
\end{graybox}
\end{column}
\begin{column}{0.48\textwidth}
\begin{graybox}[height=3.4cm]
\textbf{Image resolution}\\
768 versus 336 pixels: about a fifth of the image tokens. Details suffer first (police at the edge,
text on a placard); the overall scene suffers least.
\end{graybox}
\end{column}
\end{columns}

\vspace{0.25cm}
Also worth checking when relevant: the order of images in a multi-image prompt, the model size,
and a newer version of the same model.
\end{frame}

% ─────────────────────────────────────────────────────────────────────────────
\begin{frame}{Scaling Up: Two Routes}
\begin{columns}[T]
\begin{column}{0.48\textwidth}
\begin{greenbox}[height=4.3cm]
\textbf{Describe, then analyze}\\[2pt]
The VLM writes two sentences per image; text methods do the rest (here TF-IDF and a logistic
regression).\\[3pt]
\textcolor{darkgreen}{+} transparent: you can read what the model mentions\\
\textcolor{warnborder}{$-$} slow, and whatever isn't put into words is lost
\end{greenbox}
\end{column}
\begin{column}{0.48\textwidth}
\begin{navybox}[height=4.3cm]
\textbf{Distillation}\\[2pt]
The VLM codes a few hundred images; a linear probe on CLIP embeddings learns from those codes and
then labels the whole collection in minutes.\\[3pt]
+ fast at any scale\\
$-$ inherits the VLM's errors and adds its own
\end{navybox}
\end{column}
\end{columns}

\vspace{0.2cm}
\muted{Either way, the final classifier is validated on human codes like anything else.}
\end{frame}

% ─────────────────────────────────────────────────────────────────────────────
\begin{frame}{What to Report}
\begin{enumerate}
\setlength{\itemsep}{3pt}
\item \textbf{The model, exactly}: name, size, weights version, library version, precision,
      decoding settings
\item \textbf{The full prompt}, and how and how often answers broke the format
\item \textbf{Image handling}: resolution, cropping, anything removed
\item \textbf{Validation} per variable on a random human-coded sample: kappa or F1, and the share
      coded yes by humans and by the model; per group if you compare groups
\item \textbf{Sensitivity} to wording and resolution, and whether the findings hold
\item \textbf{Hallucination checks} for variables the model could invent
\item \textbf{If you distilled}: the VLM-labeled training set and the final classifier's validation
\end{enumerate}
\end{frame}

% ─────────────────────────────────────────────────────────────────────────────
\notebookslide{
\begin{itemize}
\setlength{\itemsep}{5pt}
\nbitem{\faComments}{Loading Qwen3-VL-2B and asking it about images}
\nbitem{\faClipboardList}{The module 04 codebook as a JSON prompt, validated variable by variable}
\nbitem{\faGhost}{Catching hallucinations with leading questions}
\nbitem{\faSlidersH}{Sensitivity to prompt wording and image resolution}
\nbitem{\faAlignLeft}{Describe, then analyze with a text classifier}
\nbitem{\faCompress}{Distilling VLM labels into a CLIP probe}
\end{itemize}

\vspace{0.35cm}
\muted{\footnotesize part3\_transformers/10\_vlms\_as\_annotators/notebooks/\\
10\_vlms\_as\_annotators\_exercises.ipynb · 6 exercises · GPU needed}
}

\end{document}
